% !TEX root = ../bb_AiRa_Kappaleet.tex

%% HARRIS: Chapter 10

% ö = \%c3\%b6
% ä = \%C3\%A4

\xdef\sectiontitle{\IIsectiontitle} %aiheuttama
\TOCrefs

%%%%%%%%%%%%%%%
\begin{frame}[label=2_7_a]%[label=2_6_TOC1]
\frametitle{\texorpdfstring{\culine[\sectiontitleulinecolor]{\sectiontitle}}{\sectiontitle} }
\label{\TOCname}

\vspace{\chapterTOCvksip}
\begin{itemize}[leftmargin=0.5cm,itemsep=3mm]
    \item[2.1]<1-> \hyperlink{2_1_a}{\IIaaSubsectiontitle}
    \item[2.2]<1-> \hyperlink{2_2_a}{\IIaSubsectiontitle}
    \item[2.3]<1-> \hyperlink{2_3_a}{\IIbSubsectiontitle}	
    \item[2.4]<1-> \hyperlink{2_4_a}{\IIcSubsectiontitle}	
    \item[2.5]<1-> \hyperlink{2_5_a}{\IIdSubsectiontitle}
    \item[2.6]<1-> \hyperlink{2_6_a}{\IIeSubsectiontitle}
    \item[2.7]<1-> \hyperlink{2_7_a}{\CDAlert[red]{\IIfSubsectiontitle}}
    \item[2.8]<1-> \hyperlink{2_8_a}{\IIgSubsectiontitle}
    \item[2.9]<1-> \hyperlink{2_9_a}{\IIhSubsectiontitle}
    \item[2.10]<1-> \hyperlink{2_10_a}{\IIiSubsectiontitle}
\end{itemize}

%\endofslide 
\end{frame}
%%%%%%%%%%%%%%%

\xdef\sectiontitle{\IIfSubsectiontitle} 

%%%%%%%%%%%%%%%
\begin{frame}
\frametitle{\texorpdfstring{\culine[\sectiontitleulinecolor]{\sectiontitle}}{\sectiontitle} }
\framesubtitle{Motivaatio puolijohdekomponentteihin}
\appear{}

\begin{itemize}[itemsep=1.8mm]
	\item Sähköisen signaalin \CDAlert[red]{päälle-}\appear{\CDAlert[red]{/poiskytkeminen}} \appear{ja \CDAlert[blue]{vahvistaminen}} \appear{ovat tärkeitä modernien sähköisten komponenttien ominaisuuksia} \appear{(viestinnässä}\appear{, kännyköissä}\appear{, tietokoneissa}\appear{, näytöissä}\appear{, valaistuksessa}\appear{, jne.)} 
	
	\item \Alert[fuchsia]{Yksinkertaiset} \CDAlert[fuchsia]{vastukset}\appear{, \CDAlert[fuchsia]{kondensaattorit}} \appear{\Alert[fuchsia]{ja} \CDAlert[fuchsia]{käämit}} \appear{\Alert[fuchsia]{eivät pysty näihin asioihin}}\appear{, täytyy hyödyntää puolijohdekomponentteja} 
	
	\item Tällaisia komponentteja ovat mm.:
	\item[] \begin{itemize} \seminormal
	\item  \CDAlert[fuchsia]{Diodi} \appear{ja \CDAlert[fuchsia]{Transistori}}
\item Metallioksidipuolijohteista tehty \appear{kanavatransistori} \appear{(MOSFET)}
\item \CDAlert[violet]{Valoa emittoiva diodi}, eli loistediodi, eli ledi \appear{(LED)}
\item \CDAlert[violet]{Ilmaisimet} \appear{(tuleva valo/elektronit synnyttävät sähkövirran)}
\item Varauskytketyt piirit \appear{(CCD:t)}
\item \CDAlert[fuchsia]{Komplementaarinen metallioksidipuolijohdeteknologia (CMOS)}
\item $\ldots$
\end{itemize}


\end{itemize}


\endofslide 
%\endofsection
\end{frame}
%%%%%%%%%%%%%%%

%%%%%%%%%%%%%%%
\begin{frame}
\frametitle{\texorpdfstring{\culine[\sectiontitleulinecolor]{\sectiontitle}}{\sectiontitle} }
\framesubtitle{Hallin efekti}
\appear{}

\begin{itemize}
	\item Puolijohteissa on kahden polarireetin varauksenkuljettajia: \appear{negatiiviset \CDAlert[blue]{elektronit}} \appear{ja positiiviset \CDAlert[red]{aukot}} 
	
	\item Tämä voidaan varmistaa kokeellisesti hyödyntämällä \CDAlert[red]{Hallin efektiä}. \appear{Virta laitetaan kulkemaan seostetun puolijohteen läpi}\appear{, joka on ulkoisessa magneettikentässä.}\appear{ Lorentzin voima poikkeuttaa varauksenkuljettajia sivusuunnassa}\appear{, kohtisuorasti sähkövirtaan nähden: }
\[
\vec{F}  \equal  \appear{q} \appear{(\vec{v}} \appear{\times \vec{B})}
\]

\end{itemize}

\xdef\loc{south}
\begin{tikzpicture}[remember picture,overlay] % anchor=south
    \node (A) [yshift=-0.3*\figYshift,xshift=0*\figXshift, anchor=\loc] at (current page.\loc) {\includegraphics[page=1,height=3.5cm]{Figures_booklet/SOM_Tipler_Solid_state_physics_figures/SOM_Tipler_SSP_figure_39_cropped.png}}; %scale=\figscale. % 8cm max height
\end{tikzpicture}
\footnote{\TiplerCR[1-]}

\endofslide 
%\endofsection
\end{frame}
%%%%%%%%%%%%%%%

%%%%%%%%%%%%%%%
\begin{frame}
\frametitle{\texorpdfstring{\culine[\sectiontitleulinecolor]{\sectiontitle}}{\sectiontitle} }
\framesubtitle{Hallin efekti}
\appear{}

\begin{itemize}
\item Koska varaustenkuljettajien polariteetti  \appear{on otettu kaavassa huomioon kahdesti}  \appear{(polariteeetti ja suunta eli nopeuden $v_d$ merkki)}\appear{, molempien polariteettien varauksenkuljettajat synnyttävät samansuuntaisen virran.}\appear{ Varauksenkuljettajien liike} \appear{(ja sen indusoima voima)} \appear{synnyttää kuitenkin lisäksi polariteetista riippuvan sähkökentän $E_{Hall}$.}

\item Indusoitunut \CDAlert[fuchsia]{Hallin kenttä $E_{Hall}$} johtaa mitattavaan potentiaalieroon 
\[
V_{\text {Hall}}  \equal  \appear{v_{d}} \appear{B} \appear{w} \appear{\,, \quad ~\text{missä}} \qquad \appear{w=\text{johteen leveys}}
\]
\end{itemize}

\xdef\loc{south}
\begin{tikzpicture}[remember picture,overlay] % anchor=south
    \node (A) [yshift=-0.3*\figYshift,xshift=0*\figXshift, anchor=\loc] at (current page.\loc) {\includegraphics[page=1,height=3.5cm]{Figures_booklet/SOM_Tipler_Solid_state_physics_figures/SOM_Tipler_SSP_figure_39_cropped.png}}; %scale=\figscale. % 8cm max height
\end{tikzpicture}
\footnote{\TiplerCR[1-]}

\endofslide 
%\endofsection
\end{frame}
%%%%%%%%%%%%%%%

%%%%%%%%%%%%%%%
\begin{frame}
\frametitle{\texorpdfstring{\culine[\sectiontitleulinecolor]{\sectiontitle}}{\sectiontitle} }
\framesubtitle{Aukot varauksenkuljettajina}
\appear{}

\begin{itemize}
	\item Aukon käsite varauksenkuljettajana on hienovarainen asia\appear{, eikä vastaa tilannetta missä syntyy elektronien liikkeen suhteen vastakkainen virta} \appear{(\CDAlert[fuchsia]{elektronivyön rakenne} vaikuttaa elektroniin ja aukkoon eri tavalla)}
	
	\item Absoluuttisessa nollapisteessä puolijohteen\appear{ valenssivyö on täysi ja johtavuusvyö taas tyhjä} 
	\implies{ Aine ei reagoi 
			ulkoiseen sähkökenttään }
\end{itemize}
	
% 6.5 cm = midway, 10 cm = 3/4 
\vspace{2.5mm}
\begin{minipage}{7.95cm}
\begin{itemize}
	\item $T$:n kasvaessa\appear{ aineessa alkaa tapahtumaan termisiä virityksiä}\appear{, jolloin elektroni virittyy hetkellisesti valenssivyöltä johta\-vuus\-vyöl\-le}\appear{ jättäen valenssivyölle aukon} 
	\implies{ \CDAlert[red]{Elektroni} voi tällöin liikkua vapaasti \\
			\Alert[red]{johtavuusvyöllä} } \vspace{-2mm}
	\implies{ \CDAlert[blue]{Aukko} voi myös vapaasti liikkua \\
			 \Alert[blue]{valenssivyöllä} }
\end{itemize}
\end{minipage}
\vspace{2.5mm}

\xdef\loc{south east}
\begin{tikzpicture}[remember picture,overlay] % anchor=south
    \node (A) [yshift=-0.25*\figYshift,xshift=\figXshift, anchor=\loc] at (current page.\loc) {\includegraphics[page=1,height=6cm]{Figures_booklet/SOM_10_Bonding_figures/SOM_Bonding_Figure_35.png}}; %scale=\figscale. % 8cm max height
\end{tikzpicture}
\footnote{\HarrisCR[1-]}

\endofslide 
%\endofsection
\end{frame}
%%%%%%%%%%%%%%%


%%%%%%%%%%%%%%%
\begin{frame}
\frametitle{\texorpdfstring{\culine[\sectiontitleulinecolor]{\sectiontitle}}{\sectiontitle} }
\framesubtitle{Aukot varauksenkuljettajina}

\begin{itemize}
	\item<1-> Aukon käsite varauksenkuljettajana on hienovarainen asia, eikä vastaa tilannetta missä syntyy elektronien liikkeen suhteen vastakkainen virta (\CDAlert[fuchsia]{elektronivyön rakenne} vaikuttaa elektroniin ja aukkoon eri tavalla)
	
	\item<1-> Absoluuttisessa nollapisteessä puolijohteen\appear{ valenssivyö on täysi ja johtavuusvyö taas tyhjä} 
	\Implies{ Aine ei reagoi 
			ulkoiseen sähkökenttään }
\end{itemize}	
% 6.5 cm = midway, 10 cm = 3/4 
\vspace{2.5mm}
\begin{minipage}{7.9cm}
\begin{itemize}
	\item Aukko ei itsessään liiku\appear{ vaan se mahdollistaa vieressä olevien elektronien liikkeen} 
	\item \Alert[fuchsia]{Toinen elektroni voi siirtyä aukon paikalle}\appear{ jättäen jälkeensä siirtyneen aukkopaikan}
	\implies{ Vaikkeivät valenssielektronit ole itsessään vapaita\appear{, tyhjä tila eli aukko}\appear{ voi vapaasti liikkua } }
\end{itemize}
\end{minipage}
\vspace{2.5mm}

\xdef\loc{south east}
\begin{tikzpicture}[remember picture,overlay] % anchor=south
    \node (A) [yshift=-0.25*\figYshift,xshift=\figXshift, anchor=\loc] at (current page.\loc) {\includegraphics[page=1,height=6cm]{Figures_booklet/SOM_10_Bonding_figures/SOM_Bonding_Figure_35.png}}; %scale=\figscale. % 8cm max height
\end{tikzpicture}
\footnote{\HarrisCR[1-]}

\endofslide 
%\endofsection
\end{frame}
%%%%%%%%%%%%%%%

%%%%%%%%%%%%%%%
\begin{frame}
\frametitle{\texorpdfstring{\culine[\sectiontitleulinecolor]{\sectiontitle}}{\sectiontitle} }
\framesubtitle{Aukot varauksenkuljettajina}
\appear{}

\begin{itemize}
	\item Huomaa energiatilojen käyttäytyminen aukkojen noustessa/laskeutuessa energiakaaviossa:
\item[] \begin{itemize}[itemsep=3mm]
	 \item Kun aukko nousee vyöllä\appear{ kokonaisenergia laskee}
	 \end{itemize}
\end{itemize}

% 6.5 cm = midway, 10 cm = 3/4 
\vspace{2.5mm}
\begin{minipage}{7.6cm}
\begin{itemize}
\item[] \begin{itemize}[itemsep=3mm]
	\item Ilmiötä voidaan verrata\\ \CDAlert[fuchsia]{vedessä nousevaan ilmakuplaan}. \appear{Kun kupla nousee}\appear{ osa vedestä laskeutuu}\appear{, johtaen systeemin ko\-ko\-nais\-e\-ner\-gian pienenemiseen.}
\item Elektronit laskeutuvat spontaanisti\appear{ mutta aukot nousevat.} \appear{ Kun aukko liikkuu valenssivyöllä}\appear{, sen energia muuttuu vastakkaisella tavalla kuin elektronin tapauksessa}
\end{itemize}
\end{itemize}
\end{minipage}
\vspace{2.5mm}


\xdef\loc{south east}
\begin{tikzpicture}[remember picture,overlay] % anchor=south
    \node (A) [yshift=-0.45*\figYshift,xshift=\figXshift, anchor=\loc] at (current page.\loc) {\includegraphics[page=1,height=6cm]{Figures_booklet/SOM_10_Bonding_figures/SOM_Bonding_Figure_36.png}}; %scale=\figscale. % 8cm max height
\end{tikzpicture}
\footnote{\HarrisCR[1-]}


\endofslide 
%\endofsection
\end{frame}
%%%%%%%%%%%%%%%

%%%%%%%%%%%%%%%%
%\begin{frame}
%\frametitle{\texorpdfstring{\culine[\sectiontitleulinecolor]{\sectiontitle}}{\sectiontitle} }
%\framesubtitle{Note on the density of states}
%
%\begin{itemize}
%	\item We saw that the density of states for free electrons in a quantum well formed a parabolic curve: $D=D(E) \propto \sqrt{E}$, increasing with increasing energy. The actual shape of $D$, depends on structure of the crystal lattice and the position of the band. In fact, it needs to be that the density is almost symmetric for the band: 
%	\item Suppose that a band is completelty filled. If one electron is removed (perhaps to the conduction band...), then it is possible to say that a hole has been created, since a vacant state is there. When an external force is applied, some elecron may be moved into the vacant state. But by doing so, the electron leaves a new hole, corresponding to a state that was previously occupied. We can say that the hole moves in a manner opposite to electron, i.e. it 'floats'. The zero energy for the holes is at the top of the band and their energy is measured downwards, i.e. $\quad\left(E-E_{\min }\right) \rightarrow\left(E_{\max }-E\right)$
%	
%\end{itemize}
%
%%\xdef\loc{east}
%%\begin{tikzpicture}[remember picture,overlay] % anchor=south
%%    \node (A) [yshift=\figYshift,xshift=\figXshift, anchor=\loc] at (current page.\loc) {\includegraphics[page=1,height=7cm]{Figures_booklet/SOM_10_Bonding_figures/SOM_Bonding_Figure_1.png}}; %scale=\figscale. % 8cm max height
%%\end{tikzpicture}
%
%\endofslide 
%%\endofsection
%\end{frame}
%%%%%%%%%%%%%%%%

%%%%%%%%%%%%%%%
\begin{frame}
\frametitle{\texorpdfstring{\culine[\sectiontitleulinecolor]{\sectiontitle}}{\sectiontitle} }
\framesubtitle{Varauksenkuljettajan efektiivinen massa}
\appear{}

\seminormal
\begin{itemize}
	\item Jos liikkuvaa hiukkasta käsitellään aallon tavoin\appear{ sen nopeutta kuvataan ryhmänopeudella$^*$:\footnoteleft{\appear[.-]{$^*$Tilanne on analoginen ultralyhyen valopulssin nopeuden kanssa. Eikä ole mitenkään triviaali asia.}}}\vspace{-4mm}
\[
v_{\text {hiukkanen}}  \equal \appear{v_{\text {ryhmä}}}  \equal \frac{d \omega}{d k}
\]
\item Ryhmänopeuden funktiota $\omega=\omega(k)$\appear{, kutsutaan "\CDAlert[fuchsia]{dispersiorelaatioksi}".} \appear{Hiukkasen kiihtyvyydeksi saadaan sen avulla:}
\[
a  \equal \frac{d \textrm{((}v_{\text {hiukkanen}} \textrm{))}}{d t}  \equal \frac{d}{d t} \LP \frac{d \omega}{d k} \,\RP  \equal \frac{d^{2} \omega}{d k^{2}} \frac{d k}{d t}
\]

\item Tarkastellaan ulkoisen voiman vaikutusta hiukkaseen.\appear{ Yksikköajassa tehty työ (teho) saadaan kertomalla voima nopeudella:}
$$
\begin{aligned}
\appear{F_{\text {ext}}} \appear{v_{\text {hiukkanen}}} & \equal \frac{d E}{d t}  \equal \frac{d(\hbar \omega)}{d t}  \equal \appear{\hbar} \frac{d \omega}{d t} \equal \appear{\hbar} \frac{d \omega}{d k} \frac{d k}{d t}  \equal \appear{\hbar v_{\text {hiukkanen}}} \frac{d k}{d t} \qquad \\
\Rightarrow \qquad \qquad \appear{F_{ext}}&  \equal \appear{\hbar} \frac{d k}{d t}
\end{aligned}
$$
\end{itemize}

\endofslide 
%\endofsection
\end{frame}
%%%%%%%%%%%%%%%

%%%%%%%%%%%%%%%
\begin{frame}
\frametitle{\texorpdfstring{\culine[\sectiontitleulinecolor]{\sectiontitle}}{\sectiontitle} }
\framesubtitle{Varauksenkuljettajan efektiivinen massa}
\appear{}

\begin{itemize}
	\item Käytetään seuraavaksi Newtonin 2:sta lakia \appear{$(F=m a)$} \appear{apuna määrittelemään varauksenkuljettajien \CDAlert[fuchsia]{efektiivinen massa}:}
\end{itemize}
% 6.5 cm = midway, 10 cm = 3/4 
%\vspace{2.5mm}
\begin{minipage}{8.2cm}
\begin{itemize}
	\item[] 
	$$
	\begin{aligned}
m^{*} & \equal \frac{F_{e x t}}{a}  \equal \frac{\hbar \frac{d k}{d t}}{\frac{d^{2} \omega}{d k^{2}} \frac{d k}{d t}}  \equal \frac{\hbar}{\frac{d^{2} \omega}{d k^{2}}}  \equal \appear{\hbar^{2} }\LP \frac{d^{2} E}{d k^{2}} \,\RP\appear{\raisebox{8pt}{$^{-1}$}} \qquad \\
\Rightarrow \qquad \frac{1}{m^{*}}&  \equal \frac{1}{\hbar^{2}} \frac{d^{2} E}{d k^{2}}
\end{aligned}
$$ %\appear{\textrm{))}^{-1}} 

\item Varmistetaan vielä että tämä määritelmä \appear{on yhteensopiva (vapaan) kappaleen energian kanssa:}
\[
 E  \equal \frac{p^{2}}{2 m_e}  \equal \frac{\hbar^{2} k^{2}}{2 m_e}
\]
\end{itemize}
\end{minipage}
\vspace{2.5mm}


\xdef\loc{south east}
\begin{tikzpicture}[remember picture,overlay] % anchor=south
    \node (A) [yshift=-0.35*\figYshift,xshift=\figXshift, anchor=\loc] at (current page.\loc) {\includegraphics[page=1,height=7cm]{Figures_booklet/SOM_10_Bonding_figures/SOM_Bonding_Figure_27.png}}; %scale=\figscale. % 8cm max height
\end{tikzpicture}
\footnote{\HarrisCR[1-]}

\endofslide 
%\endofsection
\end{frame}
%%%%%%%%%%%%%%%


%%%%%%%%%%%%%%%
\begin{frame}
\frametitle{\texorpdfstring{\culine[\sectiontitleulinecolor]{\sectiontitle}}{\sectiontitle} }
\framesubtitle{Varauksenkuljettajan efektiivinen massa}
\appear{}

% 6.5 cm = midway, 10 cm = 3/4 
%\vspace{2.5mm}
\begin{minipage}{8.2cm}
\begin{itemize}
	\item[] %\vspace{-4mm}
	$$
 E =\Frac{p^{2}}{2 m_e}=\Frac{\hbar^{2} k^{2}}{2 m_e}
$$
\item Katsotaan kuvaa 27\appear{ mikä näyttää elektro\-nien energian \textrm{((}massa $m_{e} $\textrm{))} kiinteässä aineessa} \appear{aaltoluvun $k$ funktiona} 

\item Energian toinen derivaatta $k$:suhteen joh\-taa tällöin yhtälöön
\[
\frac{d^{2} E}{d k^{2}}  \equal  \appear{\hbar^{2}} \frac{1}{m_{e}} \quad \Rightarrow \quad \appear{m_{e}} \equal \appear{\hbar^{2}} \LP \frac{d^{2} E}{d k^{2}} \,\RP\appear{\raisebox{8pt}{$^{-1}$}} \equiv \appear{m^*}
\]
\item Eli efektiivisen massan $m^*$ käsite \appear{on yhteensopiva elektronin massan kanssa}\appear{, kun elektroni käyttäytyy vapaan hiukkasen tavoin } \appear{(keskellä vyötä)}
\end{itemize}
\end{minipage}
\vspace{2.5mm}


\xdef\loc{south east}
\begin{tikzpicture}[remember picture,overlay] % anchor=south
    \node (A) [yshift=-0.35*\figYshift,xshift=\figXshift, anchor=\loc] at (current page.\loc) {\includegraphics[page=1,height=7cm]{Figures_booklet/SOM_10_Bonding_figures/SOM_Bonding_Figure_27.png}}; %scale=\figscale. % 8cm max height
\end{tikzpicture}
\footnote{\HarrisCR[1-]}

\endofslide 
%\endofsection
\end{frame}
%%%%%%%%%%%%%%%

%%%%%%%%%%%%%%%
\begin{frame}
\frametitle{\texorpdfstring{\culine[\sectiontitleulinecolor]{\sectiontitle}}{\sectiontitle} }
\framesubtitle{Varauksenkuljettajan efektiivinen massa}
\appear{}

\semismall
\begin{itemize}
%	\item[] \vspace{-2mm}
%	$$
%m_{e}=\hbar^{2} \textrm{((}\Frac{d^{2} E}{d k^{2}} \textrm{))}^{-1} \qquad \& \qquad \frac{1}{m_{e}} \equal \frac{1}{\hbar^{2}} \frac{d^{2} E}{d k^{2}}
%$$
\item $E=E(k)$:n kaareutuvuus \appear{(= toinen derivaatta)}  \appear{liittyy massan käänteisarvoon (massan käänteisluku) }
\end{itemize}

\vspace{1.5mm}

\begin{minipage}{7.8cm}
\begin{itemize}[itemsep=1.5mm]

\item Kun pinkki käyrä seuraa katkoviivaa\appear{, eli paraabelia $\appear{E} \propto \appear{k^{2}}$}\appear{, energia vastaa vapaan hiukkasen energiaa}\appear{, ja pätee $m^{*}=m_{e}$}

\item Lähellä vöiden reunoja\appear{, energia-aukon molemmilla puolilla} \appear{efektiivinen massa poikkeaa $m_{e}$:stä.} \appear{Aloitetaan arvosta $k=0$}\appear{, lähellä arvoa $k=\pi / a$ kaarevuus kasvaa suureksi} \appear{ja lähellä yläreu\-naa muuttuu jopa negatiiviseksi.} \appear{Energiaraon yläpuolella kaarevuus on suuri ja positiivinen, jne.}

\item Alueilla missä aukon energia on pieni aukon le\-vey\-teen verrattuna\appear{, \Alert[blue]{efektiivinen massa $m^{*}$ voi olla $0,01$- tai $0,1$-kertainen vapaan elektronin massaan nähden}}
\end{itemize}
\end{minipage}

\xdef\loc{south east}
\begin{tikzpicture}[remember picture,overlay] % anchor=south
    \node (A) [yshift=-0.35*\figYshift,xshift=\figXshift, anchor=\loc] at (current page.\loc) {\includegraphics[page=1,height=7cm]{Figures_booklet/SOM_10_Bonding_figures/SOM_Bonding_Figure_37.png}}; %scale=\figscale. % 8cm max height
\end{tikzpicture}
\footnote{\HarrisCR[1-]}


\endofslide 
%\endofsection
\end{frame}
%%%%%%%%%%%%%%%

%%%%%%%%%%%%%%%
\begin{frame}
\frametitle{\texorpdfstring{\culine[\sectiontitleulinecolor]{\sectiontitle}}{\sectiontitle} }
\framesubtitle{Johtavuus elektronien ja aukkojen seurauksena}
\appear{}

\seminormal

\begin{itemize}
	\item Tavallisissa johteissa (kuten kuparissa)\appear{, elektronit täyttävät keskellä kaistaa olevat tilat}\appear{, missä ne käyttäytyvät vapaan hiukkasen tavoin $m\Approx m_e$} 
	\item Puolijohteessa \appear{elektronit täyttävät tilat valenssivyön huipulle asti}\appear{, missä efektiivinen massa on negatiivinen}
	\item Aukon energia määräytyy käänteisesti \appear{(valenssi-}
\end{itemize}

% 6.5 cm = midway, 10 cm = 3/4 
%\vspace{2.5mm}
\begin{minipage}{8.5cm}
\begin{itemize}
	\item[]  ja johtavuusvyön kaarevuudet ovat aukoille vastakkaiset). \appear{Vyön huipulla oleva aukko} \appear{liikkuu ulkoisen kentän suuntaan} \appear{aukon positiivinen massan ansiosta}

\item Sekä johtavuuselektronit \appear{että valenssiaukot} \appear{toi\-mi\-vat \CDAlert[fuchsia]{varauksenkuljettajina}} \appear{ja reagoivat ul\-koi\-seen kenttään}\appear{, ja vaikuttavat sähkönjoh\-ta\-vuuteen:}\vspace{-0.5mm}
\[
\qquad \sigma  \equal \frac{\appear{n_{e}} \appear{e^{2}} \appear{\tau_{e}}}{\appear{m_{e}^{*}}}  \plus \frac{n_{h} e^{2} \tau_{h}}{m_{h}^{*}}
\]
\end{itemize}
\end{minipage}

\xdef\loc{south east}
\begin{tikzpicture}[remember picture,overlay] % anchor=south
    \node (A) [yshift=-0.35*\figYshift,xshift=\figXshift, anchor=\loc] at (current page.\loc) {\includegraphics[page=1,height=6cm]{Figures_booklet/SOM_10_Bonding_figures/SOM_Bonding_Figure_38.png}}; %scale=\figscale. % 8cm max height
\end{tikzpicture}
\footnote{\HarrisCR[1-]}

\endofslide 
%\endofsection
\end{frame}
%%%%%%%%%%%%%%%

%%%%%%%%%%%%%%%
\begin{frame}
\frametitle{\texorpdfstring{\culine[\sectiontitleulinecolor]{\sectiontitle}}{\sectiontitle} }
\framesubtitle{Varaustenkuljettajien tilatiheydet}
\appear{}

%
\begin{itemize}
	\item Muistellaan tilatiheyttä $D=D(E)$ \appear{ja miehityslukua $ \mathbb{N} = \mathbb{N} (E)_{F D}$ fermioneille:}
\[
D(E)  \equal \frac{(2 s+1) m^{3 / 2} V}{\pi^{2} \hbar^{3} \appear{\sqrt{2}}} \appear{\sqrt{E} }\qquad \appear{\&} \qquad  \appear{\mathbb{N} (E)_{F D}}  \equal \frac{1}{e^{ \textrm{((}E-E_{F} \textrm{))} / k_{B} T}+1}
\]

\item Varauksenkuljettajien lukumäärätiheyttä puolijohteessa voidaan arvioida  \appear{integroimalla $D(E) \cdot  \mathbb{N} (E)_{F D}$ sopivan energiavälin yli}
\end{itemize}

% 6.5 cm = midway, 10 cm = 3/4 
\vspace{2.5mm}
\begin{minipage}{7cm}
\begin{itemize}
	\item Olkoon $E_{c}$ ja $E_{v}$ johtavuusvyön pohjan ja valenssivyön huipun energiat

\item Tällöin energiarako on yksinkertaisesti $\appear{E_g} \equal \appear{E_{c}} \minus \appear{E_{v}}$
\end{itemize}
\end{minipage}

\xdef\loc{south east}
\begin{tikzpicture}[remember picture,overlay] % anchor=south
    \node (A) [yshift=-0.25*\figYshift,xshift=1*\figXshift, anchor=\loc] at (current page.\loc) {\includegraphics[page=1,height=4cm]{Figures_booklet/SOM_10_Bonding_figures/Tikz_SOM_Bonding_figures5}}; %scale=\figscale. % 8cm max height
\end{tikzpicture}

\endofslide 
%\endofsection
\end{frame}
%%%%%%%%%%%%%%%

%%%%%%%%%%%%%%%
\begin{frame}
\frametitle{\texorpdfstring{\culine[\sectiontitleulinecolor]{\sectiontitle}}{\sectiontitle} }
\framesubtitle{Varaustenkuljettajien tilatiheydet}
\appear{}

\begin{itemize}
	\item \CDAlert[blue]{Negatiivisille varauksenkuljettajille} $(s=1 / 2)$ kirjoitetaan $D(E) \times  \mathbb{N} (E)_{F D}$:
\[
D(E) \appear{\cdot  \mathbb{N} (E)_{F D}} \equal \fraca{2 m_{e}^{* 3 / 2} V}{\pi^{2} \hbar^{3} 2^{1/2} \appear{(E-E_{c})^{1/2}}} \appear{\cdot} \fraca{1}{e^{ \textrm{((}E-E_{F} \textrm{))} / k_{B} T}+1}
\]
\item[] Tätä voidaan jatkossa käyttää sellaisenaan

\item Mutta \CDAlert[red]{positiivisilla varauksenkuljettajilla} $(s=1 / 2)$\appear{, eli aukoilla }\appear{ valenssivyön muoto on erilainen huipulla}\appear{, jolloin energian neliöjuuri pitää korvata termillä $(E_{v}-E)^{1/2}$}

\item Lisäksi Fermin–Diracin jakauma $  \mathbb{N} (E)_{F D}=\Frac{1}{e^{ \textrm{((}E-E_{F} \textrm{))} / k_{B} T}+1}$
\appear{täytyy korvata termillä} $\appear{1}  \minus  \frac{1}{e^{ \textrm{((}E-E_{F} \textrm{))} / k_{B} T}+1}$

\item Yhdistetty todennäköisyys millä aukko ilmenee tietyllä energialla\appear{, ja todennäköisyys että tila on miehitty}\appear{  täytyy summautua yhdeksi} \appear{(\Alert[fuchsia]{elektronin täytyy olla \CDAlert[fuchsia]{jossain}})}
%\item The probability of hole occurring for certain energy when added to the probability of the state being occupied, must be unity
\end{itemize}

\endofslide 
%\endofsection
\end{frame}
%%%%%%%%%%%%%%%

%%%%%%%%%%%%%%%
\begin{frame}
\frametitle{\texorpdfstring{\culine[\sectiontitleulinecolor]{\sectiontitle}}{\sectiontitle} }
\framesubtitle{Varaustenkuljettajien tilatiheydet}
\appear{}

\begin{itemize}
	\item Saamme:
\[
1 \minus \frac{1}{e^{\Frac{E-E_{F}}{k_{B} T}}+1} \equal \frac{e^{\Frac{E-E_{F}}{k_{B} T}}}{e^{\Frac{E-E_{F}}{k_{B} T}}+1} \equal \frac{1}{1+e^{\Frac{- \textrm{((}E-E_{F} \textrm{))}}{k_{B} T}}} \equal \frac{1}{1+e^{\Frac{E_{F}-E}{k_{B} T}}} \approx \appear{e^{\Frac{- \textrm{((}E_{F}-E \textrm{))}}{k_{B} T}}}
\]

\item Jolloin positiivisille varauksenkuljettajille:
\[
D(E) \appear{\times  \mathbb{N} (E)_{F D}} \equal \frac{2 m_{h}^{* 3 / 2} V}{\pi^{2} \hbar^{3} \sqrt{2}} \appear{\sqrt{E_{v}-E}} \appear{e^{- \textrm{((}E_{F}-E \textrm{))} /k_{B} T}}
\]
\item Nyt positiivisten ja negatiivisten varauksenkuljettajien tilatiheydet voidaan laskea
\end{itemize}

\endofslide 
%\endofsection
\end{frame}
%%%%%%%%%%%%%%%

%%%%%%%%%%%%%%%
\begin{frame}
\frametitle{\texorpdfstring{\culine[\sectiontitleulinecolor]{\sectiontitle}}{\sectiontitle} }
\framesubtitle{Varaustenkuljettajien tilatiheydet}
\appear{}

\begin{itemize}
	\item Negatiivisille varauksenkuljettajille $n_{e}$ saamme:
$$
\begin{aligned} 
\appear{n_{e}}&  \equal  \frac{1}{V} \appear{\nint_{E_{c}}^{\infty}} \frac{2 m_{e}^{* 3 / 2} V }{\pi^{2} \hbar^{3} \sqrt{2}} \appear{\sqrt{E-E_{c}}} \appear{e^{\frac{- \textrm{((}E-E_{F} \textrm{))}}{k_{B} T}}} \appear{d E}  \equal \frac{2 m_{e}^{* 3 / 2} }{\pi^{2} \hbar^{3} 2^{1/2}} \appear{e^{E_F / k_{B}T}} \appear{\nint_{E_{c}}^{\infty}} \appear{\sqrt{E-E_{c}}} \appear{e^{\Frac{-E}{k_{B} T}} d E} \\ 
& \equal \frac{1}{4} \LP \frac{2 m_{e}^{*} k_{B} T}{\pi \hbar^{2}} \,\RP\appear{\raisebox{8pt}{$^{3/2}$}}  \appear{e^{- \textrm{((}E_{c}-E_{F} \textrm{))} / k_{B} T} }
\end{aligned}
$$

\item Positiivisille varauksenkuljettajille $n_{h}$ saamme:
$$
\begin{aligned} 
\appear{n_{h}}&  \equal \frac{1}{V} \appear{\nint_{-\infty}^{E_{v}}} \frac{2 m_{h}^{* 3 / 2} V}{\pi^{2} \hbar^{3} \sqrt{2}} \appear{\sqrt{E_{v}-E}} \appear{e^{\frac{- \textrm{((}E_{F}-E \textrm{))}}{k_{B} T}} \appear{d E}}  \equal \frac{2 m_{e}^{* 3 / 2} }{\pi^{2} \hbar^{3} 2^{1/2}} \appear{e^{-\Frac{E_F}{k_BT}}}  \appear{\nint_{-\infty}^{E_{v}}} \appear{\sqrt{E_{v}-E}} \appear{e^{\frac{E}{k_{B} T}} \appear{d E}  }\\ 
& \equal \frac{1}{4} \appear{\Bigg(} \frac{2 m_{h}^{*} k_{B} T}{\pi \hbar^{2}} \appear{\Bigg)^{3 / 2}} \appear{e^{- \textrm{((}E_{F}-E_{v} \textrm{))} / k_{B} T}} \\ 
\end{aligned}
$$
%$$
%\begin{aligned} 
%\appear{n_{h}}&  \equal \frac{1}{V} \appear{\nint_{-\infty}^{E_{v}}} \frac{2 m_{h}^{* 3 / 2} V}{\pi^{2} \hbar^{3} \sqrt{2}} \appear{\sqrt{E_{v}-E}} \appear{e^{\frac{- \textrm{((}E_{F}-E \textrm{))}}{k_{B} T}} \appear{d E}}  \equal \frac{2 m_{e}^{* 3 / 2} V}{\pi^{2} \hbar^{3} 2^{1/2}} \appear{e^{-\Frac{E_F}{k_BT}}}  \appear{\nint_{-\infty}^{E_{v}}} \appear{\sqrt{E_{v}-E}} \appear{e^{\frac{E}{k_{B} T}} \appear{d E}  }\\ 
%& \equal \frac{1}{4} \appear{\Bigg(} \frac{2 m_{h}^{*} k_{B} T}{\pi \hbar^{2}} \appear{\Bigg)^{3 / 2}} \appear{e^{- \textrm{((}E_{F}-E_{v} \textrm{))} / k_{B} T}} \\ 
%\end{aligned}
%$$
\item Nämä tulokset lasketaan kotitehtävissä !! 
\end{itemize}

\endofslide 
%\endofsection
\end{frame}
%%%%%%%%%%%%%%%

%%%%%%%%%%%%%%%
\begin{frame}
\frametitle{\texorpdfstring{\culine[\sectiontitleulinecolor]{\sectiontitle}}{\sectiontitle} }
\framesubtitle{Massatoiminnan laki: \quad $n_e n_h =n_i^2$}
\appear{}

\begin{itemize}
	\item Eli tulokseksi on saatu:
\[
n_{e} \appear{n_{h}}  \equal \frac{1}{2} \LP \frac{k_{B} T}{\pi \hbar^{2}} \,\RP\appear{\raisebox{8pt}{$^{\!3}$}} \appear{\textrm{((}m_{e}^{*} m_{h}^{*} \textrm{))}^{3 / 2}} \appear{e^{-\Delta E / k_{B} T}} \appear{\,, \qquad \Delta E}  \equal  \appear{E_g}  \equal  \appear{E_c - E_v}
\]

\item Huomaa että \CDAlert[blue]{itseispuolijohteilla} \appear{negatiivisten ja positiivisten varauksenkuljettajien tilatiheydet ovat yhtäsuuret:}
\[
n_{e}  \equal  \appear{n_{h}}  \equal  \frac{1}{4} \LP \frac{2 k_{B} T}{\pi \hbar^{2}} \,\RP\appear{\raisebox{8pt}{$^{\!3/2}$}} \appear{\textrm{((}m_{e}^{*} m_{h}^{*} \textrm{))}^{3 / 4}} \appear{e^{-\Delta E / 2 k_{B} T}} \equiv \appear{n_i^2}
\]
\item[] missä $n_i$ on varauksenkuljettajien tiheys\appear{, jonka avulla voidaan ilmaista erittäin hyödyllinen  \CDAlert[red]{massatoiminnan laki}:} 
\[
n_{e} \appear{n_{h}}  \equal \appear{n_i} \qquad \appear{\textrm{(pätee myös \Alert[red]{seostetuille puolijohteille}!)}}
\]
\item Huomaa että varauksenkuljettajien lukumäärät riippuvat sekä lämpötilasta  \appear{että energiaraosta $\Delta E$}\appear{, eli riippuvat myös materiaalista }
\end{itemize}

\endofslide 
%\endofsection
\end{frame}
%%%%%%%%%%%%%%%

%%%%%%%%%%%%%%%
\begin{frame}
\frametitle{\texorpdfstring{\culine[\sectiontitleulinecolor]{\sectiontitle}}{\sectiontitle} }
\appear{}

\framesubtitle{Itseispuolijohteet}
% On intrinsic semiconductors

\seminormal
\begin{itemize}[itemsep=1.6mm]
%	\item Previously, we discussed materials which are defined as \Alert[fuchsia]{intrinsic semiconductors}\appear{, undoped elements} \appear{(or compounds)} \appear{with a band gap smaller than $2$\;eV} \appear{(the choice of $2$\;eV is somewhat arbitrary)} 
	
	\item Tähän asti olemme käsitelleet \CDAlert[red]{alkuainepuolijohteita} \appear{(eli Si}\appear{, Ge)}\appear{, tai \CDAlert[blue]{yhdistelmäpuolijohteita}} \appear{(esim. SiC}\appear{, Si$_{1-x}$Ge$_{x}$}\appear{, GaAs)}\appear{, mitkä kaikkia ovat olleet \CDAlert[fuchsia]{itseispuolijohteita}}\appear{, eli seostamattomia puolijohteita}  \appear{joiden energiarako on pienempi kuin $2$\;eV} \appear{($2$\;eV on vähän mielivaltainen raja)} 
	
	\item Niin pienellä energiaraolla\appear{ suhteellisen matalakin lämpötila riittää aiheuttamaan johtavuusvyön merkittävän miehityksen}
	
	\end{itemize}
% 6.5 cm = midway, 10 cm = 3/4 
\vspace{2.5mm}
\begin{minipage}{5.8cm}
\begin{itemize}[itemsep=1.6mm]
%\item[] significant occupation of the conduction band
	\item \CDAlert[red]{Pii} \appear{ja \CDAlert[blue]{germanium},} 
	\appear{ovat kaksi yleistä itseispuolijohdetta} 
	\appear{joiden energiarako on \CDAlert[red]{$1,14$\;eV}} \appear{ja \CDAlert[blue]{$0,67$\;eV}, vastaavasti}

\item Molemmat ovat kovalenttisia aineita\appear{, joissa atomit jakavat kaikki 4 valenssielektroniaan}\appear{, piille $\mathrm{n}=3$ kuorella}\appear{, germaniumille $(Z=32)$} \appear{$n=4$ kuorella}
\end{itemize}
\end{minipage}
\vspace{2.5mm}


\xdef\loc{south east}
\begin{tikzpicture}[remember picture,overlay] % anchor=south
    \node (A) [yshift=-0.25*\figYshift,xshift=\figXshift, anchor=\loc] at (current page.\loc) {\includegraphics[page=1,height=5cm]{Figures_booklet/SOM_Tipler_Solid_state_physics_figures/SOM_Tipler_SSP_figure_36_cropped.png}}; %scale=\figscale. % 8cm max height
\end{tikzpicture}
\footnote{\TiplerCR[1-]}

\endofslide 
%\endofsection
\end{frame}
%%%%%%%%%%%%%%%

%%%%%%%%%%%%%%%
\begin{frame}
\frametitle{\texorpdfstring{\culine[\sectiontitleulinecolor]{\sectiontitle}}{\sectiontitle} }
\framesubtitle{Itseispuolijohteet}
\appear{}

\seminormal
\begin{itemize}[itemsep=2.2mm]
	\item Kahden Si atomin kovalenttinen sidos, jossa elektronitiheys on korkea atomien välissä \appear{näkyy kuvassa b, joka on muodostettu käyttämällä röntgensirontaa}

\item Arseeni \appear{$(As, Z=33)$} \appear{on amorfisessa muodossaa myös alkuainepuolijohde jonka energiarako on $1,2-1,4$\,eV.} \appear{Arseenilla on kuitenkin  \CDAlert[red]{viisi valenssielektronia}} \appear{(eli se on \CDAlert[red]{pentavalentti})} \appear{ $n=4$ kuorella.} \appear{Neljä viidestä valenssielektronista osallistuu kovalenttisten sidosten muodostumiseen}\appear{, mutta \CDAlert[red]{viides elektroni} on vain \CDAlert[red]{löyhästi sidoksissa} atomiin}

\item Arseenia käytetään useammin (käyttöä halutaan vähentää) muodostamaan \CDAlert[blue]{III-V puolijohteita} (GaAs), tai \CDAlert[red]{seosaineena seostetuissa puolijohteissa} 

\item Arseenilla seostettaessa 5:s elektroni asettuu vain hieman johtavuusvyön alapuolella olevalle energiatilalle\appear{, josta se voidaan helposti virittää johtavuusvyölle.} \appear{Tämä viides elektroni ja arseenin ioninen ydin muodostavat \CDAlert[tunired]{vetyatomimaisen systeemin}}\appear{, jonka energiatilat ovat laskettavissa vastaavasti}

%\item In pure materials, the number of electrons thermally excited to the conduction band \appear{necessarily equals the number of holes left behind in the valence band}
\end{itemize}

\endofslide 
%\endofsection
\end{frame}
%%%%%%%%%%%%%%%

%%%%%%%%%%%%%%%
\begin{frame}
\frametitle{\texorpdfstring{\culine[\sectiontitleulinecolor]{\sectiontitle}}{\sectiontitle} }
\framesubtitle{Seostetut puolijohteet}
%Doping of semiconductors $-$ Extrinsic semiconductors
\appear{}

\begin{itemize}
	\item Seostamalla\appear{, eli lisäämällä pieni määrä epäpuhtausalkuaineita}\appear{, voidaan muodostaa \CDAlert[fuchsia]{seostettuja puolijohteita}} 
	
	\item Näillä \Alert[fuchsia]{on suurempi määrä joko johtavuuselektroneja tai valenssiaukkoja} \appear{sekä}\appear{ paljon suurempi varauksenkuljettajien konsentraatio kuin itseispuolijohteissa}

\item Esimerkiksi kovalenttinen Si \appear{(tai Ge hila)}\appear{, \CDAlert[red]{jota on seostettu pentavalentilla arseenilla}}
\end{itemize}

\appear{\xdef\loc{south}
\begin{tikzpicture}[remember picture,overlay] % anchor=south
    \node (A) [yshift=-0.25*\figYshift,xshift=0*\figXshift, anchor=\loc] at (current page.\loc) {\includegraphics[page=2,height=4cm]{Figures_booklet/SOM_10_Bonding_figures/Tikz_SOM_Bonding_figures6_12_fixed}}; %scale=\figscale. % 8cm max height
\end{tikzpicture}
}

\endofslide 
%\endofsection
\end{frame}
%%%%%%%%%%%%%%%

%%%%%%%%%%%%%%%
\begin{frame}
\frametitle{\texorpdfstring{\culine[\sectiontitleulinecolor]{\sectiontitle}}{\sectiontitle} }
\appear{}

\begin{itemize}
	\item Jatketaan arseenin viidennen elektronin tarkastelua
	
	\item \CDAlert[tunired]{Bohrin mallia} \appear{voidaan käyttää energiatilojen arviointiin:}
	\[
E_{n}  \equal  \minus \frac{1}{2} \LP \frac{k e^{2}}{\hbar} \,\RP \frac{m_{e}}{n^{2}} \appear{\times} \frac{m^{*}}{m_{e}}  \frac{1}{\kappa^{2}} \equal  \minus \frac{1}{2} \LP \frac{k e^{2}}{\hbar} \,\RP \frac{m^{*}}{\kappa^{2}}  \frac{1}{n^{2}}
\]
\appear{missä $\mathrm{m}^{*}$ on efektiivinen massa}\appear{, $\kappa$ on dielektrinen vakio}\appear{, ja Bohrin ratojen keskimääräiset säteet ovat} 
\[
 < r_{n}>\,  \equal \appear{a_{0}} \appear{n^{2}}  \frac{m_{e}}{m^{*}}  \appear{\kappa} \appear{\,, \qquad a_{0}} \appear{= \text{Bohrin säde}~0,0529~\mathrm{nm} }
 \] 

\item Tämä energia on laskettu suhteessa johtavuusvyön minimienergiaan $E_c$\appear{, koska arseeni on epäpuhtaus piihilassa }
\implies{ yksi ylimääräinen elektroni voi vapaasti liikkua ja johtaa sähköä 
 }
\item $E_{\infty}=0$ on energialla $E_c$\appear{, muut energiat $E_{n}$ ovat sen lähellä sen alapuolella} \appear{energiaraon sisällä }

\end{itemize}

\endofslide 
%\endofsection
\end{frame}
%%%%%%%%%%%%%%%

%%%%%%%%%%%%%%%
\begin{frame}
\frametitle{\texorpdfstring{\culine[\sectiontitleulinecolor]{\sectiontitle}}{\sectiontitle} }
\framesubtitle{n-tyypin seostetut puolijohteet}
\appear{}

\seminormal
\begin{itemize}[itemsep=1.65mm]
	
\item Jos oletetaan että $m^* \Approx 0,2 m_{e}$ \appear{(piin $m^{*}$)} \appear{niin tila $E_{1}=-0,020$\;eV on vain hieman johtavuusvyön alapuolella} \appear{($E_{c}=0$\;eV)}  \appear{eli reilusti heikompi kuin $-13,6$\;eV (= vedyn perustila)}

\item Myös $\appear{<r_{1}>=3,1 \mathrm{~nm}} \appear{=60 \times}$\appear{vedyn perustilan elektronin keskimääräinen säde}

\item Ylimääräiset elektronit voivat siis helposti virittyä johtavuusvyölle\appear{, koska \CDAlert[red]{ionisaatioenergia on lähellä huoneenlämmön $k T$ arvoa}}

\item Vetymäisiä tasoja johtavuusvyön alapuolella kutsutaan \CDAlert[red]{luovuttajatiloiksi}\appear{, koska ne \Alert[red]{'luovuttavat' elektronin johtavuusvyölle}} \appear{jättämättä aukkoa valenssivyölle}

\item Tämän tyypin puolijohteita kutsutaan \CDAlert[red]{n-tyypin puolijohteiksi}

\item Seostetun aineen johtavuuteen voidaan vaikuttaa muuttamalla epäpuhtauksien konsentraatiota 
\implies{ Jo muutaman miljoonasosan (parts per million) \CDAlert[red]{(ppm) määrä epäpuhtauksia} vaikuttaa huomattavasti johtavuuteen }
\end{itemize}

\endofslide 
%\endofsection
\end{frame}
%%%%%%%%%%%%%%%

%%%%%%%%%%%%%%%%
%\begin{frame}
%\frametitle{\texorpdfstring{\culine[\sectiontitleulinecolor]{\sectiontitle}}{\sectiontitle} }
%
%%n-tyoe extrinsic semiconductor
%\begin{itemize}
%	\item Thus, the extra electrons can be easily excited to the conduction band, since the ionization energy is close to $k T$ of room temperature
%
%\item The hydrogen-like levels just below the conduction band are called donor levels, since they 'donate' electrons to the conduction band without leaving holes in the valence band.
%
%\item This type of impurity semiconductor is called $\underline{n}$-type semiconductor
%
%\item The conductivity of the doped material can be controlled by controlling the amount of impurity added
%
%\item Even few parts per million (ppm) of impurity
%\end{itemize}
%
%%\xdef\loc{east}
%%\begin{tikzpicture}[remember picture,overlay] % anchor=south
%%    \node (A) [yshift=\figYshift,xshift=\figXshift, anchor=\loc] at (current page.\loc) {\includegraphics[page=1,height=7cm]{Figures_booklet/SOM_10_Bonding_figures/SOM_Bonding_Figure_1.png}}; %scale=\figscale. % 8cm max height
%%\end{tikzpicture}
%
%\endofslide 
%%\endofsection
%\end{frame}
%%%%%%%%%%%%%%%%

%%%%%%%%%%%%%%%
\begin{frame}
\frametitle{\texorpdfstring{\culine[\sectiontitleulinecolor]{\sectiontitle}}{\sectiontitle} }
\framesubtitle{n-tyypin seostetut puolijohteet}
\appear{}

\begin{itemize}
	\item Itseispuolijohteen valenssivyö on täysin täynnä \implies{ Jokainen epäpuhtausatomi \appear{lisää ylimääräisen elektronin johtavuusvyölle} }

\item Jokainen epäpuhtaus-ioni on myös positiivisesti varautunut \appear{(kokonaisvaraus säilyy neutraalina)}

\item Kokonaisuutena syntyy luovuttajatila\appear{ joka on heti johtavuusvyön alla}\appear{, tyypillisesti ${\sim} 0,05$\;eV matalammalla}
\end{itemize}
\vspace{2mm}
\appear{\xdef\loc{south east}%
\begin{tikzpicture}[remember picture,overlay]% anchor=south
    \node (A) [yshift=-0.25*\figYshift,xshift=\figXshift, anchor=\loc] at (current page.\loc) {\includegraphics[page=1,width=7cm]{Figures_booklet/SOM_10_Bonding_figures/SOM_Bonding_Figure_39.png}};%scale=\figscale. % 8cm max height
\end{tikzpicture}\footnote{\HarrisCR}%
}
\begin{minipage}{6.5cm}
\begin{itemize} 
	\item Luovuttajatilan elektronit \appear{on helppo virittää johtavuusvyölle} \appear{ jo\\ matalissa lämpötiloissa} 
\end{itemize}

\end{minipage}

\endofslide 
%\endofsection
\end{frame}
%%%%%%%%%%%%%%%

%%%%%%%%%%%%%%%
\begin{frame}
\frametitle{\texorpdfstring{\culine[\sectiontitleulinecolor]{\sectiontitle}}{\sectiontitle} }
\framesubtitle{p-tyypin seostetut puolijohteet}
\appear{}

\seminormal
%p-type extrinsic semiconductor
\begin{itemize}[itemsep=1.6mm]
	\item On olemassa myös ns. \CDAlert[blue]{$\mathbf{p}$-tyypin puolijohteita}
	
	\item Esim. Si:stä koostuvaan puolijohteeseen voidaan lisätä \CDAlert[blue]{trivalenttinen} \Alert[blue]{alumiini- \appear{(tai gallium-)}}\appear{atomi}\appear{, jolla on kolme valenssielektronia}
	
	\item Alumiini mielellään vastaanottaa yhden elektronin sitä ympäröiviltä Si atomeilta muodostaakseen neljä täyttä kovalenttista sidosta
	\implies{  Näitä tiloja kutsutaan '\CDAlert[blue]{vastaanottajatiloiksi}' tai akseptoritiloiksi }

	\item Akseptoritilat ottavat mielellään elektroneita vastaan ympäröivien atomien valenssivöiltä\appear{, mihin riittää jo huoneenlämmön $k_B T$ määrä energiaa}
\end{itemize}

\appear{\xdef\loc{south}
\begin{tikzpicture}[remember picture,overlay] % anchor=south
    \node (A) [yshift=-0.05*\figYshift,xshift=0*\figXshift, anchor=\loc] at (current page.\loc) {\includegraphics[page=3,height=4cm]{Figures_booklet/SOM_10_Bonding_figures/Tikz_SOM_Bonding_figures6_12_fixed}}; %scale=\figscale. % 8cm max height
\end{tikzpicture}
}

\endofslide 
%\endofsection
\end{frame}
%%%%%%%%%%%%%%%

%%%%%%%%%%%%%%%
\begin{frame}
\frametitle{\texorpdfstring{\culine[\sectiontitleulinecolor]{\sectiontitle}}{\sectiontitle} }
\framesubtitle{p-tyypin seostetut puolijohteet}
\appear{}

\begin{itemize}
	\item Kun elektroni virittyy vastaanottajatilaan\appear{, valenssivyölle muodostuu aukko}
	\implies{ Puolijohteeseen muodostuu positiivisia varauksenkuljettajia }
	
	\item Yhteenvetona seostusaineen\appear{, eli lisätyn epäpuhtauden} \appear{tyyppi}\appear{ vaikuttaa tuleeko itseispuolijohteesta} \appear{\Alert[fuchsia]{joko} \CDAlert[red]{n-tyypin}} \appear{\Alert[fuchsia]{vai} \CDAlert[blue]{p-tyypin} \Alert[fuchsia]{seostettu puolijohde}}
\end{itemize}

% 6.5 cm = midway, 10 cm = 3/4 
\vspace{2.5mm}
\begin{minipage}{6.5cm}
	\begin{itemize}	
	\item \Alert[red]{Johtavuuksia voidaan suuresti kasvattaa}! 
	
	\item Vielä tärkeämpää on\appear{ että muodostamalla n- ja p-tyypin puolijoh\-de\-alueita} \appear{on mahdollista valmistaa ihmeellisiä  \Alert[fuchsia]{sähköisiä komponentteja} (puolijohdekomponentteja)! }
\end{itemize}
\end{minipage}

\xdef\loc{south east}
\begin{tikzpicture}[remember picture,overlay] % anchor=south
    \node (A) [yshift=-0.35*\figYshift,xshift=1*\figXshift, anchor=\loc] at (current page.\loc) {\includegraphics[page=1,width=7cm]{Figures_booklet/SOM_10_Bonding_figures/SOM_Bonding_Figure_40.png}}; %scale=\figscale. % 8cm max height
\end{tikzpicture}
\footnote{\HarrisCR[1-]}

%\endofslide 
\endofsection
\end{frame}
%%%%%%%%%%%%%%%

